\documentclass[11pt]{article}
\usepackage{amsmath,amssymb,amsthm,booktabs,geometry,graphicx}
\usepackage[utf8]{inputenc}
\usepackage[T1]{fontenc}
\usepackage[french]{babel}
\usepackage[round,authoryear]{natbib}
\usepackage[colorlinks=true,linkcolor=black,citecolor=black,urlcolor=blue!50!black]{hyperref}
\geometry{margin=1in}
\setlength{\emergencystretch}{2em}

\newtheorem{proposition}{Proposition}
\newtheorem{corollary}{Corollaire}
\newtheorem{assumption}{Hypothèse}
\theoremstyle{definition}
\newtheorem{postulate}{Postulat}
\newtheorem{definition}{Définition}
\theoremstyle{remark}
\newtheorem*{remark}{Remarque}
\renewcommand{\proofname}{Démonstration}

\newcommand{\JJ}{\mathcal{J}}
\newcommand{\DD}{\mathcal{D}}
\newcommand{\Dun}{\mathcal{D}_1}
\newcommand{\KL}[2]{D_{KL}\!\left(#1\,\Vert\,#2\right)}

\title{Compléments théoriques\\[4pt]
\large Fondements informationnels d'un indicateur entropique\\de stress systémique : ce que le papier principal laisse de côté}
\author{Annexe au papier \og Sauts d'entropie et relaxation\\de l'information sur les marchés financiers \fg{}}
\date{Octobre 2026}

\begin{document}
\maketitle

\begin{abstract}
\noindent Le papier principal est recentré sur un seul objet destiné aux
banques centrales : le canal de dépendance $-\tfrac12\log\det R$ d'un panel
multi-marchés, sa décomposition exacte entre blocs et à l'intérieur des blocs,
et l'audit en nats des scénarios de stress publiés. Ce document rassemble tout
ce que ce recentrage retire du corps du texte, pour servir d'annexe ou être
fourni sur demande. On y trouve le pont qui mène du maximum d'entropie en
temps-volume au mouvement brownien géométrique, la décomposition de l'entropie
en trois canaux et l'indice structurel $\JJ=\KL{p}{\prod_i q_i}$, la dynamique
de saut et diffusion (un théorème, deux postulats, la structure
d'Ornstein--Uhlenbeck imposée par le théorème de Doob), le programme empirique
H1--H8 avec son verdict face aux témoins calibrés, et la géométrie complète des
scénarios en boule de Kullback--Leibler. Deux résultats sont nouveaux par
rapport à la version arXiv~v1. Le premier est une identité de décomposition par
blocs du canal de dépendance, qui relie ce document au papier principal. Le
second montre que le relèvement d'un scénario de portefeuille à la dimension
pleine qui minimise la divergence ambiante coûte exactement le rayon $\eta$, et
qu'il s'interprète actif par actif (choc porté par la covariance de chaque actif
avec le portefeuille), alors que la dilatation uniforme par défaut coûte $N$~fois
plus en variance. Ce second point corrige une remarque de la v1.
\end{abstract}

\tableofcontents

\section*{Avant-propos : ce que contient ce document, et pourquoi}
\addcontentsline{toc}{section}{Avant-propos}

\paragraph{L'articulation avec le papier principal.} L'évaluation d'octobre 2026
a conclu qu'un seul résultat de la version arXiv~v1 franchit les deux témoins
calibrés sur les deux panels : l'écart de régime du canal de dépendance
$H^{dep}=\tfrac12\log\det R$. La phase~A l'a confirmé sous trois formes (brute,
dévolatilisée à la Forbes--Rigobon, en rangs) et a montré que le ratio
d'absorption de \citet{kritzman2011}, mesure de référence de la même famille,
échoue au témoin par blocs. Le papier principal se reconstruit donc autour d'un
\og CISS entropique \fg{} au sens de \citet{hollo2012} : canal de dépendance d'un
panel multi-marchés, décomposition intra/inter-blocs, validation à la manière
des banques centrales, et audit des scénarios publiés en nats.

Tout le reste sort du corps du texte. Ce n'est pas parce que c'est faux : les
résultats mathématiques restent exacts. C'est parce que ces éléments ne sont
pas nécessaires à l'indicateur, ou parce qu'ils ne survivent pas aux témoins et
ne peuvent donc pas fonder un outil de surveillance. Les rassembler ici a deux
usages. Un rapporteur qui demande d'où vient le canal de dépendance, et pourquoi
c'est lui et pas l'indice structurel complet, trouve la réponse. Un lecteur qui
veut reprendre la théorie de la dynamique de l'information trouve la chaîne
entière, avec le statut exact de chaque maillon.

\paragraph{Statut de chaque énoncé.} Le document distingue trois catégories et
les signale au fil du texte : les \emph{théorèmes} (exacts sous leurs
hypothèses), les \emph{postulats} (hypothèses économiques substantielles,
testables), et les \emph{constats empiriques} (avec le témoin qu'ils
franchissent ou non). Le tableau~\ref{tab:statut} en donne la carte.

\begin{table}[htbp]
\centering
\small
\begin{tabular}{@{}p{5.0cm}p{3.0cm}p{4.6cm}l@{}}
\toprule
Élément & Statut & Verdict empirique & Section \\
\midrule
Maximum d'entropie gaussien & théorème & --- & \ref{sec:bridge} \\
Pont volume (MDH, brownien subordonné) & argument sous hypothèses & non testé ici & \ref{sec:bridge} \\
Décomposition en trois canaux & identité exacte & --- & \ref{sec:channels} \\
Indice structurel $\JJ$ comme divergence & identité exacte & étendue franchit les deux témoins & \ref{sec:channels} \\
Sauts d'entropie vers le bas & théorème (en moyenne) & non estimable directement & \ref{sec:dynamics} \\
Relaxation (P1) & postulat & demi-vie non identifiée & \ref{sec:dynamics} \\
Contrainte unilatérale (P2) & postulat & non établie (H5) & \ref{sec:dynamics} \\
Couplage queues--corrélation & conséquence & non identifié (H6) & \ref{sec:dynamics} \\
Canal de dépendance sous stress & constat & \textbf{franchit les deux témoins} & \ref{sec:program}, \ref{sec:blocks} \\
Boule KL, $1$~nat $=\sqrt2\,\sigma$ & théorème & --- & \ref{sec:stress} \\
Relèvement minimal (rang un) & théorème (nouveau) & --- & \ref{sec:stress} \\
Calibration du rayon sur le flux & règle & trop douce face à l'histoire & \ref{sec:stress} \\
Décomposition par blocs & identité (nouvelle ici) & objet du papier principal & \ref{sec:blocks} \\
\bottomrule
\end{tabular}
\caption{Carte des énoncés du document et de leur statut.}
\label{tab:statut}
\end{table}

\paragraph{Origine des textes.} Les sections~\ref{sec:bridge} à
\ref{sec:dynamics} et les annexes~\ref{app:units} à \ref{app:mi} reprennent, avec
des retouches de renvois, le texte de la version française de référence
(\texttt{paper/main\_fr.tex}). Les sections~\ref{sec:program} à
\ref{sec:open} sont rédigées pour ce document. Tous les chiffres cités
proviennent du même code (\texttt{python -m src.run\_empirical}, témoins à
deux cents réplications, phase~A) ; aucun n'a été recalculé pour cette annexe.

